% Kansrekenen. Caravaggio: kaarten en vals spel, het kansspel waaruit de kansrekening ontstond.
% Motief: de 36 even waarschijnlijke uitkomsten van twee dobbelstenen, gestapeld per som.
\def\omslagbestand{omslag-valsspelers.jpg}
\def\omslagbron{Caravaggio, \textit{De valsspelers}, ca. 1594 (publiek domein)}
\def\omslagrand{10.5 + 0.35*sin(\x*360/14)}
\def\omslagmotief{%
  \randlijn
  \foreach \som/\aantal in {2/1,3/2,4/3,5/4,6/5,7/6,8/5,9/4,10/3,11/2,12/1}{%
    \pgfmathsetmacro{\plek}{13.9+(\som-2)*.62}%
    \foreach \k in {1,...,\aantal}{%
      \ifnum\som=7
        \fill[fel,draw=accent,line width=.6pt] (\plek,{8.75+(\k-1)*.62}) circle (2.2mm);
      \else
        \fill[omslagstip] (\plek,{8.75+(\k-1)*.62}) circle (2.2mm);
      \fi
    }
  }
  \foreach \som in {2,7,12}{%
    \pgfmathsetmacro{\plek}{13.9+(\som-2)*.62}%
    \node[text=white,font=\small,inner sep=1pt,anchor=north] at (\plek,8.4) {$\som$};
  }
}
